\documentclass[UTF8,11pt]{article}

\usepackage{ctex}
\usepackage[a4paper,margin=1in]{geometry}
\usepackage{amsmath,amssymb,amsthm,mathtools}
\usepackage{algorithm}
\usepackage{algpseudocode}
\usepackage{xcolor}
\usepackage{booktabs}
\usepackage{hyperref}
\usepackage{enumitem}
\usepackage{microtype}

\hypersetup{colorlinks=true,linkcolor=blue,citecolor=blue,urlcolor=blue}

\newcommand{\todo}[1]{\textcolor{red}{\textbf{[待完成：}#1\textbf{]}}}
\newcommand{\RN}{R_{q,N}}
\newcommand{\Rn}{R_{q,n}}
\newcommand{\Split}{\mathsf{Split}}
\newcommand{\Merge}{\mathsf{Merge}}
\newcommand{\RSDown}{\mathsf{RSDown}}
\newcommand{\RSUp}{\mathsf{RSUp}}
\newcommand{\CtwoS}{\mathsf{C2S}}
\newcommand{\StwoC}{\mathsf{S2C}}
\newcommand{\EvalMod}{\mathsf{EvalMod}}
\newcommand{\logSlots}{\mathsf{logSlots}}
\newcommand{\CtwoSScaling}{\mathsf{C2SScaling}}
\newcommand{\StwoCScaling}{\mathsf{S2CScaling}}
\newcommand{\core}{\mathsf{core}}
\newcommand{\BTS}{\mathsf{BTS}}
\newcommand{\KS}{\mathsf{KS}}
\newcommand{\Enc}{\mathsf{Enc}}
\newcommand{\Dec}{\mathsf{Dec}}
\newcommand{\polyR}[1]{\mathbb{Z}_{#1}[X]/(X^N+1)}
\newcommand{\ct}{\mathbf{ct}}
\newcommand{\repmsg}{\mathrel{\vDash_{\Delta}}}

\newcommand{\ZZ}{\mathbb Z}
\newcommand{\RR}{\mathbb R}
\newcommand{\CC}{\mathbb C}
\newcommand{\II}{\mathrm i}

\newcommand{\SplitKeep}{\mathsf{Split}^{\mathrm{keepY}}_{n,k}}
\newcommand{\SplitNorm}{\mathsf{Split}^{\mathrm{norm}}_{n,k}}
\newcommand{\MergeKeep}{\mathsf{Merge}^{\mathrm{keepY}}_{n,k}}
\newcommand{\MergeNorm}{\mathsf{Merge}^{\mathrm{norm}}_{n,k}}
\newcommand{\etak}{\eta_k}
\newcommand{\etaTwo}{\eta_{2k}}
\newcommand{\repd}{\mathrel{\dashv\vdash}}

\newcommand{\Round}{\operatorname{Round}}


\newcommand{\Rbig}{\mathcal{R}_N^{\RR}}
\newcommand{\Rint}{\mathcal{R}_N^{\ZZ}}
\newcommand{\Ssmall}{\mathcal{S}_n}

\newcommand{\rep}{\mathrel{\vDash_{\Delta}}}

\newcommand{\Z}{\mathbb{Z}}
\newcommand{\C}{\mathbb{C}}
\newcommand{\R}{\mathbb{R}}
\newcommand{\Q}{\mathbb{Q}}


\newcommand{\SN}{\mathcal S_N}
\newcommand{\Sn}{\mathcal S_n}


\newcommand{\ctP}{\mathbf{ct}_P}
\newcommand{\ctminus}{\mathbf{ct}_{-}}

\newcommand{\cte}{\mathbf{ct}_{\mathrm{even}}^{\uparrow}}
\newcommand{\ctoy}{\mathbf{ct}_{\mathrm{odd},Y}^{\uparrow}}
\newcommand{\cto}{\mathbf{ct}_{\mathrm{odd}}^{\uparrow}}

\newcommand{\Pe}{P_{\mathrm{even}}}
\newcommand{\Po}{P_{\mathrm{odd}}}


\newcommand{\Ecd}{\mathsf{Ecd}}
\newcommand{\Dcd}{\mathsf{Dcd}}

\newcommand{\Auto}[1]{\mathsf{Auto}_{#1}}

\newcommand{\Rot}{\mathsf{Rot}}
\newcommand{\Conj}{\mathsf{Conj}}
\newcommand{\Add}{\mathsf{Add}}
\newcommand{\Sub}{\mathsf{Sub}}
\newcommand{\MulPt}[1]{\mathsf{MulPt}_{#1}}

\newcommand{\RS}{\mathsf{Rescale}}


\newcommand{\Snh}{\mathcal S_{n/2}}

\newcommand{\Comp}{\mathsf{Compress}}



\newcommand{\Level}{\mathsf{Level}}
\newcommand{\ksnoise}{\eta_{\mathrm{ks}}}



\newcommand{\Fold}{\mathsf{Fold}}

\newcommand{\Compress}{\mathsf{Compress}}
\newcommand{\Expand}{\mathsf{Expand}}


\newcommand{\Tr}{\mathsf{Tr}}
\newcommand{\Prd}{\mathsf{Pr}}
\newcommand{\Diag}{\mathsf{Diag}}



\newcommand{\round}[1]{\left\lfloor #1 \right\rceil}
\newcommand{\norm}[1]{\left\lVert #1 \right\rVert}

\theoremstyle{definition}
\newtheorem{definition}{定义}
\newtheorem{remark}{注}
\theoremstyle{plain}
\newtheorem{lemma}{引理}
\newtheorem{theorem}{定理}
\newtheorem{proposition}{命题}
\newtheorem{example}{例如}

\title{基于子环 Ring Switching 的 CKKS Bootstrapping 分解\\
\large 初代成果汇总版本}
\author{梁朝阳}
\date{\today}

\begin{document}
\maketitle

\begin{abstract}
CKKS bootstrapping 的主要开销来自线性变换和模约减核心。本文研究一个大环 CKKS bootstrapping 是否能够被分解为若干个小环上的 leaf bootstrapping。给定 $N=kn$，我们将 HERMES-style 子环 ring switching 与明文多项式的 coefficient/module decomposition 结合起来。本文的主要理论主张是：传统 CKKS bootstrapping core
\[
\StwoC_N\circ\EvalMod_N\circ\CtwoS_N
\]
与该分解兼容。换言之，split 之后可以在每个 leaf 上独立执行
\[
\StwoC_n\circ\EvalMod_n\circ\CtwoS_n,
\]
然后再 merge 回大环。该构造不需要在 split 之后立即恢复原始的大环 slots。本文还讨论 coefficient/slot 双视角下的 split 语义、一个更一般的两步骤拆分解释、实现诊断、参数选择，以及 reduced-modulus leaf bootstrapping 这一开放问题。

\todo{在安全估计、key-size 表格和 reduced-$Q$ 实验完成之后，替换为最终摘要。}
\end{abstract}


\tableofcontents

\newpage

\section{引言}

前面可能需要写点大同行看的


\subsection{研究动机}
CKKS bootstrapping 通过同态方式计算一个类似解密的电路，从而刷新密文。在传统视角下，ModRaise 或 ModUp 之后，明文多项式中会出现形如 $m+q_0 I$ 的项。随后 bootstrapping core 将 coefficient 信息搬到 slots 中，逐坐标执行近似 modular reduction，再将结果搬回 coefficient 表示。我们将该核心写成
\[
\BTS_N^{\mathrm{core}}=\StwoC_N\circ\EvalMod_N\circ\CtwoS_N.
\]
其中昂贵的不只是 $\EvalMod$ 中的标量近似，也包括大环上的线性变换 $\CtwoS_N$ 和 $\StwoC_N$。因此，本文的核心问题是：当 $N=kn$ 时，能否将 $R_{q,N}$ 上的一次 bootstrapping 分解为 $R_{q,n}$ 上的若干次更小的 bootstrapping？

本文的答案是：利用 subring ring switching，将一个大环 ciphertext 转换为 $k$ 个 leaf ciphertexts，在每个 leaf 上执行 bootstrapping core，然后合并输出。整体形态为
\[
\RSUp_{n\to N}\circ\left(\bigoplus_{r=0}^{k-1}\BTS_n\right)\circ\RSDown_{N\to n}.
\]
这里的 $\RSDown$ 和 $\RSUp$ 是 ciphertext-level ring switching 操作，包含 key switching，并不是单纯的 coefficient rearrangement。代数层面的 coefficient decomposition 很便宜；而密文层面的转换需要 switching keys，并会带来额外噪声。

\subsection{核心观察}
设 $P(X)\in R_N$ 且 $N=kn$。coefficient/module decomposition 将其写成
\[
P(X)=\sum_{r=0}^{k-1}X^rP_r(X^k).
\]
关键观察不是这个 split 保留了大环 slots；事实上它并不保留。真正重要的是：这个 coefficient split 与 $\CtwoS\to\EvalMod\to\StwoC$ 的结构兼容。更具体地说，split 之后执行 leaf $\CtwoS_n$，等价于先执行 top-level $\CtwoS_N$，再按照 coefficient residue classes 对坐标分组，差别只在一个 layout permutation。由于 $\EvalMod$ 是逐坐标作用，它可以在这些分组上独立执行。最后，leaf $\StwoC_n$ 加上 merge 可以恢复 top-level $\StwoC_N$ 的语义。

因此，大环 core 可以替换为以下 leaf-wise core：
\[
\Split_k\to\bigoplus_r\CtwoS_n\to\bigoplus_r\EvalMod_n\to\bigoplus_r\StwoC_n\to\Merge_k.
\]
如果从 slot-preserving split 的角度出发，某种 slot-lifting 或 slot-recovery transform 看起来会自然出现。本文的构造不将其作为算法步骤，而是直接在 coefficient/module layout 下工作。

\subsection{本文贡献}
本文计划形成三个贡献。

\paragraph{贡献一：CKKS bootstrapping 的 coefficient-split factorization。} 我们证明传统 CKKS bootstrapping core 可以在不需要额外线性变换的情况下通过 coefficient/module decomposition 分解为 leaf bootstrapping core。更具体地，对于 $N=kn$，有
\[
\StwoC_N\circ\EvalMod_N\circ\CtwoS_N
=\Merge_k\circ\left(\bigoplus_{r=0}^{k-1}\StwoC_n\circ\EvalMod_n\circ\CtwoS_n\right)\circ\Split_k,
\]
其中需要考虑 implementation layout conventions 和 approximation error。这个贡献合并了 subring-factored bootstrapping construction 和 split-compatible core theorem。它同时说明：split 之后不需要先恢复原始 big-ring slots。

\paragraph{贡献二：coefficient/slot 双视角和两步骤拆分解释。} 我们详细解释 coefficient split 与 slot-preserving split 的区别。对于 $k=2$，coefficient split
$P(X)=A(X^2)+XB(X^2)$ 将一对 big-ring slot values $P(\alpha),P(-\alpha)$ 映射为局部线性组合 $A(\theta)$ 和 $B(\theta)$，而不是直接保留原始 slot values。我们还描述一个更一般的 two-step interpretation：第一步执行 coefficient/module decomposition；第二步，如果应用确实需要某种特定 slot 语义，则再执行额外的 coordinate transform。这个 two-step view 更一般，也有理论意义，但因为会引入额外 linear-transform cost，它不是本文的主算法路径。

\paragraph{贡献三：实现、参数分析与 reduced-modulus 探索。} 我们在 Lattigo 上实现并诊断该构造，包括 layout alignment、implementation normalization、direct-vs-factored error、output level、timing 和 security-aware 参数选择。我们还分析 reduced-$Q$ 问题：ring switching 之后，leaf bootstrapping 是否可以使用更短的 modulus chain。当前证据表明，在 coefficient/RNS 层面存在代数调度自由，但现成 Lattigo preprocessing 不能通过简单重排直接下放到 leaves。因此，reduced-$Q$ 应被视为一个新的 leaf-parameter design problem，而不是 direct path 的中间态等价问题。

\todo{在最终确定 reduced-$Q$ study 能否作为主文实验之后，精炼三条贡献。}

\subsection{与已有工作的关系}

这里成正文的时候需要梳理一下行文。


\paragraph{传统 ring switching。} Ring switching 将大环 ciphertext 转换为一个或多个小环 ciphertext。在 packed setting 中，已有工作通常需要在代数分解之后执行额外的 linear transformation，以恢复或重新分配原始 slot values。本文的不同点是：我们不在 split 后立即恢复 big-ring slots。相反，我们证明 leaf $\CtwoS$ 自然产生 bootstrapping core 所需的 coordinate grouping。

\paragraph{HERMES。} HERMES 是 ring packing 和 transciphering 框架，目标是把许多 LWE-format ciphertexts packing 成一个 RLWE/CKKS ciphertext。它通过 low-modulus base ring packing、ring switching、HalfBTS 和 MLWE 技术降低时间和 key material。本文的输入输出不同：输入本来就是一个 CKKS/RLWE ciphertext，目标是刷新它自身。我们使用 HERMES-style ring switching 作为工具，但不声称 ring switching 本身是本文的新贡献。

\paragraph{Subring secret encapsulation。} Subring secret encapsulation 在 bootstrapping 前将 secret 切换到 subring dense secret。它的收益来自更小的 subring-secret Hamming weight、降低 EvalMod 或 digit-removal 成本，以及在 $\CtwoS/\StwoC$ 中使用 hoisted key switching。本文方法则是真正将 ciphertext 分解为 leaf ciphertexts，并在每个 leaf 上执行 bootstrapping core。简言之，subring secret encapsulation 是 subring-secret acceleration；本文是 full CKKS bootstrapping 的 subring-ciphertext decomposition。

\paragraph{PaCo 及相关新 bootstrapping 设计。} PaCo 通过 partial CoeffToSlot、blind rotations、circle-group modular additions 和 alternative ring structures 重新设计 CKKS bootstrapping。本文不重写 CKKS decryption equation，也不替代传统 EvalMod 路径。本文保留传统的 $\CtwoS\to\EvalMod\to\StwoC$ core，只改变该 core 的执行粒度。

\paragraph{传统 CKKS bootstrapping 优化。} 许多工作优化传统 bootstrapping pipeline 内部的 linear transforms、EvalMod approximation、precision 或参数选择。本文方法与这些工作正交：我们不是只优化某一个矩阵或近似多项式，而是将整个 core factor 成多个更小的 cores。

\todo{补完整 bibliographic citations，并加入一张简洁的 related-work comparison table。}
\section{Preliminaries}
\label{sec:preliminaries}

本节只介绍后文需要的 CKKS bootstrapping 背景、若干实现原语、ring switching 抽象和安全参数。我们不在这里完整介绍 CKKS 方案，也不详细展开 coefficient/slot split 的语义；后者将在 Section~\ref{sec:coeff-slot-semantics} 中专门讨论。本文主要关注的问题是：给定一个需要 bootstrapping 的 CKKS 密文，能否通过子环 ring switching 将其 bootstrapping core 分解为多个 leaf bootstrapping cores。

\subsection{CKKS ciphertext 语义与核心参数}
\label{subsec:ckks-ciphertext-params}

令
\[
R_{q,N}=\mathbb Z_q[X]/(X^N+1),
\]
其中 \(N\) 通常为二的幂。一个二元 CKKS ciphertext 写作 \(\ct=(c_0,c_1)\in R_{q,N}^2\)。在 secret \(s\in R_{q,N}\) 下，它满足形如
\[
\Dec_s(\ct)=c_0+c_1s=m+e\pmod q
\]
的解密关系，其中 \(m\) 是经过 CKKS 编码后的 plaintext polynomial，\(e\) 是小误差。不同实现可能采用 \(c_0-c_1s\) 或 \(c_0+c_1s\) 的符号约定；本文统一使用 \(c_0+c_1s\)。如果某个密文在 scale \(\Delta\) 下表示消息多项式 \(P(X)\)，我们有时写作 \(\ct\vDash_{\Delta,N}P(X)\)。这个记号只表达密文的消息语义，不表示 \(P(X)\) 可见。

CKKS 同时维护 modulus chain、scale 和 level。通常 modulus chain 以 RNS primes 的乘积表示：
\[
Q=\prod_i q_i.
\]
每次乘法和 rescale 会消耗 modulus chain 中的一部分 primes。scale \(\Delta\) 决定近似编码精度；常见实现会选择 \(2^{35}\)、\(2^{40}\) 或类似大小的 scale。bootstrapping 需要额外的 level 来执行 \(\CtwoS\)、\(\EvalMod\)、\(\StwoC\) 以及内部的 rotations、key switching 和 rescale。

除了 ciphertext modulus \(Q\)，CKKS 的 key switching、rotation keys、relinearization keys、bootstrapping keys 和本文的 ring switching keys 通常还使用 auxiliary modulus \(P\)。因此，在安全分析中不能只看输出 ciphertext 当前剩余的 \(Q\)，还需要考虑公开 evaluation keys 中可能出现的最大模数。本文用
\[
q_{\max}\approx QP
\]
作为保守的工程估计。更精确的安全声明应对 top ciphertext、leaf bootstrapping keys、split/merge switching keys 分别跑 lattice estimator。

\begin{remark}[参数 \(N,Q,P\) 的角色]
\(N\) 同时是容量参数、安全参数和性能成本参数。增大 \(N\) 可以容纳更多 slots，也可以承受更大的 \(Q\) 或 \(QP\)，但会增加 NTT、rotation、key switching、密钥和内存成本。增大 \(Q\) 或 \(P\) 可以支持更深的计算或更精确的 key switching，但会降低 RLWE 安全性，因此通常需要更大的 \(N\)。
\end{remark}

\todo{最终正文需要固定所有实现符号：\(N\) 是 ring degree；full complex slots 为 \(N/2\)；Lattigo 中的 LogSlots、half-slot convention 和本文 full embedding 记号的对应关系需要在实现章节说明。}

\subsection{CKKS bootstrapping 的 coefficient view 与 slot view}
\label{subsec:ckks-bts-two-views}

CKKS bootstrapping 的目标是刷新一个 modulus 已经接近耗尽的 ciphertext。设输入 ciphertext 在低模数 \(q_0\) 下表示 \(m(X)\)。ModRaise 或 ModUp 将它重新提升到更大的 modulus chain 中，但这一步会把消息关系变成类似
\[
m(X)+q_0 I(X)
\]
的形式，其中 \(I(X)\) 是某个 integer polynomial。bootstrapping 的任务是同态去掉 \(q_0I(X)\)，从而得到新的、具有较高 level 的 \(m(X)\) 的加密。

从 coefficient view 看，bootstrapping 处理的是 plaintext polynomial 的 coefficient-like 信息。粗略地说，\(\CtwoS\) 将这些待处理的 coefficient values 转入 slots；\(\EvalMod\) 在 slots 中对每个坐标执行近似 modular reduction；\(\StwoC\) 再把结果带回 coefficient representation。于是标准 core 可以写作
\[
\BTS_N^{\core}
=
\StwoC_N\circ\EvalMod_N\circ\CtwoS_N.
\]
本文的理论部分主要讨论这个 core，而不把 ScaleDown、ModUp、rescale schedule 和实现 normalization 混入主定理。

从 slot view 看，\(\EvalMod\) 是逐坐标函数：它在每个 slot 上独立近似 \(x\mapsto x\bmod q_0\) 或相应的 scaled modular reduction。正是这个逐坐标性，使得只要 \(\CtwoS\) 后的坐标能够按 residue class 分组，\(\EvalMod\) 就可以在 leaf 上独立执行。本文主定理的核心并不是重新设计 \(\EvalMod\)，而是证明
\[
\Split_k
\longrightarrow
\bigoplus_r \CtwoS_n
\longrightarrow
\bigoplus_r \EvalMod_n
\longrightarrow
\bigoplus_r \StwoC_n
\longrightarrow
\Merge_k
\]
在理想消息语义下等价于 top-level 的 \(\CtwoS_N\to\EvalMod_N\to\StwoC_N\)。

需要注意，本文的 \(\CtwoS\) 和 \(\StwoC\) 是 logical bootstrapping transforms。实际实现中，它们会被分解为许多 baby-step/giant-step rotations、diagonal linear transforms、rescale 和 key switching；还可能涉及 real/imag branches 或 two-ciphertext layout。这些实现细节会影响常数、level、normalization 和 key size，但不改变本节中使用的抽象结构。

\begin{remark}[为什么强调两种视角]
Coefficient split 在系数层面最自然；\(\EvalMod\) 在 slot 层面最自然。本文的证明正是连接这两个视角：split 后接 leaf \(\CtwoS_n\) 并不恢复原始 big-ring slots，而是进入与 top \(\CtwoS_N\) 输出一致的 coefficient-packed 子坐标。
\end{remark}

\subsection{同态实现中常用的基础原语}
\label{subsec:implementation-primitives}

后文会抽象使用 linear transform、ring switching 和 bootstrapping evaluator。这里先列出几个常见同态原语，以便读者理解实现层面的成本来源。

\paragraph{Automorphism.}
在 negacyclic ring \(R_{q,N}=\mathbb Z_q[X]/(X^N+1)\) 中，automorphism 通常由奇数 \(a\in(\mathbb Z/2N\mathbb Z)^*\) 指定：

\[
{\mathsf{Aut}}_{a}:\ X\mapsto X^a.
\]

在 CKKS 中，automorphism 实现 slot rotation、conjugation 或某些结构化变换，但密文执行 \({\mathsf{Aut}}_{a}\) 后通常还需要 key switching，将 secret 从 \(s(X^a)\) 切回目标 secret \(s(X)\)。例如在 \(R_{q,2n}\) 中，\({\mathsf{Aut}}_{1+2n}\) 对应 \(X\mapsto X^{1+2n}=-X\)。这个操作可以用于构造 even/odd projection 的同态版本，例如 \(F(X)\) 与 \(F(-X)\) 的组合。

\paragraph{Plaintext multiplication.}
\({\mathsf{MulPt}}_{g}(\ct)\) 表示用明文多项式 \(g(X)\) 乘密文所表示的消息。若 \(\ct\vDash P(X)\)，则理想语义下 \({\mathsf{MulPt}}_{g}(\ct)\vDash g(X)P(X)\)。在一些 branch normalization 中，需要去掉某个 \(X^t\) 因子。由于在 \(R_N\) 中 \(X^N=-1\)，有 \(X^{-t}=-X^{N-t}\)。因此实现中可能出现 \({\mathsf{MulPt}}_{-X^{N-t}}\)，它表示乘以 \(X^{-t}\)。例如若奇分支以 \(X B(X^2)\) 的形式出现，则可用 \({\mathsf{MulPt}}_{-X^{N-1}}\) 将其归一化为 \(B(X^2)\) 的语义。

\paragraph{Homomorphic linear transform.}
\(\CtwoS\)、\(\StwoC\)、slot-preserving split 中的 \(A\pm hB\)、以及一般 slot selection \(q=F_n^{-1}TF_Np\)，都可以看作 homomorphic linear transforms。实现中它们通常由 rotations、plaintext diagonal multiplications、rescale 和 key switching 组成。因此，即使某个变换在明文线性代数中很自然，在密文层面也可能代价昂贵。

\paragraph{Key switching.}
Key switching 将一个 secret 下的 ciphertext 转换到另一个 secret 下，同时保持消息语义不变并加入小的 switching error。普通 rotation、relinearization、bootstrapping key、以及本文的 ring switching 都依赖 key switching。本文后文将把这些误差统一纳入 correctness 和 parameter analysis，而在主理论等式中暂时忽略。

\begin{example}[用 \({\mathsf{Aut}}_{1+2n}\) 提取偶部]
设 \(F(X)=A(X^2)+XB(X^2)\in R_{2n}\)。因为 \({\mathsf{Aut}}_{1+2n}\) 实现 \(X\mapsto -X\)，有 \({\mathsf{Aut}}_{1+2n}(F)=F(-X)=A(X^2)-XB(X^2)\)。因此
\[
\frac{F(X)+F(-X)}{2}=A(X^2),
\qquad
\frac{F(X)-F(-X)}{2}=XB(X^2).
\]
如果希望得到 \(B(X^2)\)，还需要乘以 \(X^{-1}=-X^{2n-1}\)。密文层面，这对应 automorphism、加减法和一次 plaintext multiplication。这个例子说明：即使 even/odd split 的明文公式很简单，同态实现仍然包含 rotations/key switching 和明文乘法。
\end{example}

\subsection{Ring switching 与 coefficient/module decomposition}
\label{subsec:ring-switching-prelim}

令 \(N=kn\)，并将 \(R_{q,n}\) 视为由 \(Y=X^k\) 诱导的子环。明文层面上，任意 \(P(X)\in R_{q,N}\) 都有唯一 decomposition：
\[
P(X)=\sum_{r=0}^{k-1}X^rP_r(X^k),
\qquad P_r(Y)\in R_{q,n}.
\]
这给出了 coefficient/module split：大环作为小环上的 rank-\(k\) module，可以按 \(1,X,\ldots,X^{k-1}\) 分解。系数层面，它就是按 index modulo \(k\) 的 residue classes 分组。

但是在 ciphertext level，不能简单地重排 \(c_0,c_1\) 的系数后就得到小环 ciphertext。原因是 CKKS ciphertext 的消息通过 \(c_0+c_1s\) 隐含在 secret-dependent 解密关系中。若要得到真正的小环 ciphertext，必须把 secret 语义也切换到小环。本文把这种 lower-degree conversion 抽象写作
\[
\RSDown_{N\to n}(\ct_N)
=
(\ct_{n,0},\ldots,\ct_{n,k-1}),
\]
其理想消息语义为：若 \(\ct_N\vDash_{\Delta,N}P(X)\)，且 \(\Split_k(P)=(P_0,\ldots,P_{k-1})\)，则 \(\ct_{n,r}\vDash_{\Delta,n}P_r(Y)\)。

反向地，upper-degree conversion 将 \(k\) 个小环 ciphertexts 合并回大环 ciphertext。若 \(\ct_{n,r}\vDash_{\Delta,n}P_r(Y)\)，则理想的 \(\RSUp_{n\to N}\) 输出满足
\[
\RSUp_{n\to N}(\ct_{n,0},\ldots,\ct_{n,k-1})
\vDash_{\Delta,N}
\sum_{r=0}^{k-1}X^rP_r(X^k).
\]
概念上，可以把 \(\RSDown\) 看成 “key switching + split”，把 \(\RSUp\) 看成 “merge + key switching”。具体顺序依赖实现。例如某些实现会先构造 module ciphertext，再进行 switching；也有实现把 decomposition 和 switching 组织成一个整体 evaluator。本文理论上只需要它们满足上述消息语义。

\begin{remark}[ring switching 与 Section~\ref{sec:coeff-slot-semantics} 的关系]
Section~\ref{sec:coeff-slot-semantics} 将详细解释 \(\Split_k(P)\) 在 coefficient view 和 slot view 下分别意味着什么。本小节只强调 ciphertext-level 事实：\(\RSDown\) 和 \(\RSUp\) 是带 key-switching noise 的密文转换，不是免费明文重排。
\end{remark}

\subsection{Coefficient split 与 slot-preserving split 的最小区别}
\label{subsec:minimal-coeff-vs-slot}

为了避免与 Section~\ref{sec:coeff-slot-semantics} 重复，这里只给出最小区分。对于 \(k=2\)，coefficient split 写作
\[
P(X)=A(X^2)+XB(X^2).
\]
若 \(\theta=\alpha^2\)，则一对 big-ring slots 满足
\[
P(\alpha)=A(\theta)+\alpha B(\theta),
\qquad
P(-\alpha)=A(\theta)-\alpha B(\theta).
\]
因此 \(A(\theta)\) 和 \(B(\theta)\) 不是原始 slots \(P(\alpha),P(-\alpha)\) 的直接选择，而是它们的局部线性组合。

如果目标是 slot-preserving split，即构造小环消息 \(Q^+(Y),Q^-(Y)\)，使得 \(Q^+(\theta)=P(\alpha)\) 且 \(Q^-(\theta)=P(-\alpha)\)，则需要额外的 selector \(h(Y)\)，满足 \(h(\theta)=\alpha\)。于是
\[
Q^+(Y)=A(Y)+h(Y)B(Y),
\qquad
Q^-(Y)=A(Y)-h(Y)B(Y).
\]
在密文层面，这意味着 coefficient split 之后还要执行额外的 homomorphic linear transform。本文的 bootstrapping 分解不采用该 slot-preserving 路线；我们使用 coefficient/module split，并证明它与 \(\CtwoS\)--\(\EvalMod\)--\(\StwoC\) core 兼容。

\subsection{RLWE 安全性与参数选择}
\label{subsec:security-prelim}

CKKS 的底层安全性来自 RLWE。粗略地，攻击难度主要受 ring degree \(N\)、最大公开模数 \(q_{\max}\)、secret distribution、error distribution 和公开样本数量影响。在工程参数选择中，最重要的直觉是：
\[
N\uparrow \Rightarrow \text{安全预算增加},
\qquad
q_{\max}\uparrow \Rightarrow \text{安全性下降}.
\]
这里的 \(q_{\max}\) 不应只理解为当前 ciphertext 剩余的 modulus。对于带 evaluation keys 的 CKKS bootstrapping 系统，公开材料可能包含 rotation keys、relinearization keys、bootstrapping keys、ring switching keys 等，它们常常在 \(QP\) 上生成。因此保守估计时应取
\[
q_{\max}
=
\max\{
Q_{\mathrm{ct}},
Q_{\mathrm{eval}},
QP_{\mathrm{switch}},
QP_{\mathrm{bts}},
QP_{\mathrm{rs}}
\}.
\]
在初步工程分析中，可以先用 \(q_{\max}\approx QP\) 估计。

本文的子环分解还带来一个额外安全点：如果 leaf ciphertexts 或 leaf evaluation keys 真正在 \(R_{q,n}\) 上存在，那么它们的安全性必须按 leaf ring degree \(n\) 来估计，不能直接继承 top ring degree \(N\)。因此，整体方案安全性应写为各组件安全性的最小值：
\[
\lambda_{\mathsf{scheme}}
=
\min\{
\lambda_{\mathrm{top}},
\lambda_{\mathrm{leaf}},
\lambda_{\mathrm{RSDown}},
\lambda_{\mathrm{RSUp}},
\lambda_{\mathrm{BTSkeys}}
\}.
\]
如果某个 leaf key 或 switching key 的 RLWE 参数只有较低安全级别，则整个方案也只能声称该较低级别。

\begin{example}[为什么 leaf 安全要单独看]
假设 top ring 取 \(\log N=17\)，并使用 \(k=2\)，则 leaf ring 为 \(\log n=16\)。如果 bootstrapping 和 switching keys 的最大公开模数约为 \(\log_2(QP)\approx 1520\)，那么需要检查的是 \(n=2^{16}\) 是否能承受这一 \(q_{\max}\)，而不是只检查 top \(N=2^{17}\)。如果改用 \(k=4\) 且 top 仍为 \(\log N=16\)，则 leaf 只有 \(\log n=14\)，同样的 \(QP\) 会远超常见 128-bit 参考预算；这种设置更适合定位为 correctness prototype，而不是安全实例。
\end{example}

\todo{最终论文需要给出 lattice estimator 结果，并分别列出 top ciphertext/key、leaf bootstrapping keys、RSDown/RSUp keys 的 \(N\)、\(q_{\max}\)、secret distribution 和 estimated security bits。}

\section{子环分解的 Coefficient/Slot 语义}
\label{sec:coeff-slot-semantics}

本节解释本文所用子环分解的语义。需要首先强调的是：在同态算法中，我们实际持有和操作的对象始终是密文 \(\ct\)，而不是明文多项式 \(P(X)\)、系数向量 \(p\)，也不是 slot 向量。明文多项式、系数向量和 slot 向量只用于描述密文解密后的消息语义。因此，本节关于 coefficient split、slot split 和 slot-preserving split 的讨论，都应理解为密文所表示消息的语义，而不是明文可见操作。

本节采用 full complex canonical embedding 来说明主要关系。实际 CKKS 实现通常只使用半槽，并包含 conjugation、bit-reversal、\(5^i\) root ordering 等 layout convention；这些差异在本文理论层面都可以吸收到固定的 layout permutation 中。实现层面的 layout 和 normalization 将在实验章节单独讨论。

\subsection{密文表示语义：我们实际操作的是 \(\ct\)}
\label{subsec:ciphertext-representation-semantics}

设当前 CKKS scale 为 \(\Delta\)。若一个大环密文 \(\ct_N=(c_0,c_1)\in\mathcal R_{N,q}^2\) 在私钥 \(s_N\) 下满足
\[
c_0+c_1s_N=\Round(\Delta P(X))+e \pmod q,
\]
其中 \(P(X)\in\mathcal R_N\) 是被表示的消息多项式，\(e\) 是 CKKS 噪声项，则记作
\[
\ct_N\vDash_{\Delta,N} P(X).
\]
这里 \(\vDash_{\Delta,N}\) 表示“在大环 \(\mathcal R_N\) 和 scale \(\Delta\) 下表示”。类似地，若小环密文 \(\ct_n\in\mathcal R_{n,q}^2\) 在小环私钥 \(s_n\) 下表示 \(Q(Y)\in\mathcal R_n\)，则记作 \(\ct_n\vDash_{\Delta,n}Q(Y)\)。

如果从 CKKS 编码视角出发，一个 slot 向量 \(\mathbf z\) 会先通过编码映射变成一个消息多项式。例如在抽象记号下可写为 \(\operatorname{Ecd}(\mathbf z)=\Round(\Delta\cdot\tau^{-1}(\mathbf z))\)。所以当我们写 \(\ct\vDash_{\Delta,N}P(X)\) 时，也可以理解为该密文加密了某个 slot 向量经过 CKKS 编码后得到的多项式消息。本文后续使用 \(P(X)\)、\(Q(Y)\) 来描述消息语义，是为了避免在每个公式中都展开编码和取整。

\begin{remark}[消息语义与密文操作]
当我们写 \(\Split_k(P)=(P_0,\ldots,P_{k-1})\) 时，这只是明文语义层面的分解。真正的密文级操作应当是：给定 \(\ct_N\vDash_{\Delta,N}P(X)\)，通过 ring switching 或 key switching 产生 leaf ciphertexts \(\ct_{n,0},\ldots,\ct_{n,k-1}\)，并使得 \(\ct_{n,r}\vDash_{\Delta,n}P_r(Y)\)。这个转换不是免费重排，而是会引入 switching noise 和相应的 switching keys。
\end{remark}

\subsection{Coefficient/module split：密文下的目标语义}
\label{subsec:coefficient-module-split}

令 \(N=kn\)，并令 \(\mathcal R_N=\mathbb C[X]/(X^N+1)\)、\(\mathcal R_n=\mathbb C[Y]/(Y^n+1)\)。通过代换 \(Y=X^k\)，任意 \(P(X)\in\mathcal R_N\) 都可以唯一写成
\[
P(X)=\sum_{r=0}^{k-1}X^rP_r(X^k),\qquad P_r(Y)\in\mathcal R_n.
\]
如果 \(P(X)=\sum_{i=0}^{N-1}a_iX^i\)，则 \(P_r(Y)=\sum_{j=0}^{n-1}a_{r+jk}Y^j\)。也就是说，第 \(r\) 个小环多项式 \(P_r\) 收集的是大环消息多项式中所有下标同余于 \(r\pmod k\) 的系数。

\begin{definition}[Coefficient/module split]
明文层面的 coefficient/module split 定义为
\[
\Split_k(P)=(P_0,\ldots,P_{k-1}),
\qquad
P(X)=\sum_{r=0}^{k-1}X^rP_r(X^k).
\]
对应的 merge map 定义为
\[
\Merge_k(P_0,\ldots,P_{k-1})
=
\sum_{r=0}^{k-1}X^rP_r(X^k).
\]
\end{definition}

把这个定义放回密文语义中，coefficient split 的目标是：给定 \(\ct_N\vDash_{\Delta,N}P(X)\)，构造 \(k\) 个小环密文 \(\ct_{n,0},\ldots,\ct_{n,k-1}\)，使得
\[
\ct_{n,r}\vDash_{\Delta,n}P_r(Y),\qquad r=0,\ldots,k-1.
\]
这才是本文所谓“把大环密文拆到小环”的含义。输入密文仍然只是一对大环多项式 \((c_0,c_1)\)，输出应当是 \(k\) 个新的小环密文，并且每个输出密文在小环私钥语义下解密为对应的 \(P_r(Y)\)。

这里不能把 coefficient split 误解为“对密文分量 \(c_0,c_1\) 直接取 residue class”。原因是 \(c_0,c_1\) 本身并不等于消息 \(P\) 的系数；消息只通过解密方程 \(c_0+c_1s_N\approx \Delta P\) 隐含在密文中。如果只对 \(c_0,c_1\) 做公开系数抽取，通常得到的是绑定到大环 secret 的 module relation，而不是普通小环私钥下的 CKKS ciphertext。因此，真正的 ciphertext-level split 必须包含 ring switching 或 key switching。本文后文用 \(\operatorname{RSDown}_{N\to n}\) 表示该密文级操作，其理想语义是
\[
\operatorname{RSDown}_{N\to n}(\ct_N)
=
(\ct_{n,0},\ldots,\ct_{n,k-1}),
\qquad
\ct_{n,r}\vDash_{\Delta,n}P_r(Y).
\]

为了描述 \(P_r\) 在系数层面的结构，令 \(p=(a_0,\ldots,a_{N-1})^T\) 为 \(P\) 的系数向量。定义 residue-class 选择矩阵 \(S_r:\mathbb C^N\to\mathbb C^n\) 为
\[
S_rp=(a_r,a_{r+k},a_{r+2k},\ldots,a_{r+(n-1)k})^T.
\]
于是 \(P_r\) 的系数向量就是 \(S_rp\)。把所有 \(S_r\) 合在一起，记 \(S_kp=(S_0p,\ldots,S_{k-1}p)\)。在消息系数语义下，\(\Split_k\) 对应 \(p\mapsto S_kp\)；但在密文层面，要实现这一语义仍然需要 \(\operatorname{RSDown}_{N\to n}\)。

\begin{example}[\(k=2\) 时的二维 coefficient split]
\label{ex:k2-coeff-split}

考虑 \(N=2n\)。任意 \(P(X)\in\mathcal R_N\) 都可以唯一写成
\[
P(X)=A(X^2)+XB(X^2),
\]
其中 \(A(Y),B(Y)\in\mathcal R_n\)。若 \(P(X)=\sum_{i=0}^{2n-1}a_iX^i\)，则
\[
A(Y)=\sum_{j=0}^{n-1}a_{2j}Y^j,\qquad
B(Y)=\sum_{j=0}^{n-1}a_{2j+1}Y^j.
\]
因此，二维 coefficient split 在消息语义上就是 \(P\mapsto(A,B)\)。

在密文语义下，这句话应理解为：若输入大环密文满足 \(\ct_N\vDash_{\Delta,N}P(X)\)，那么理想的 ring-switch-down 输出应是两个小环密文 \(\ct_A,\ct_B\)，使得
\[
\ct_A\vDash_{\Delta,n}A(Y),\qquad
\ct_B\vDash_{\Delta,n}B(Y).
\]
也就是说，split 后我们希望得到的不是两个可见明文 \(A,B\)，而是两个分别表示 \(A,B\) 的 leaf ciphertexts。

该二维分解对加法是逐分量的。若 \(P(X)=A(X^2)+XB(X^2)\) 且 \(Q(X)=C(X^2)+XD(X^2)\)，则
\[
P(X)+Q(X)=(A+C)(X^2)+X(B+D)(X^2).
\]
但是该分解对乘法不是逐分量的。直接展开可得
\[
\begin{aligned}
P(X)Q(X)
&=(A(X^2)+XB(X^2))(C(X^2)+XD(X^2))\\
&=(AC)(X^2)+X(AD+BC)(X^2)+X^2(BD)(X^2).
\end{aligned}
\]
令 \(Y=X^2\)，上式可以重新写成
\[
P(X)Q(X)=\bigl(AC+YBD\bigr)(X^2)+X\bigl(AD+BC\bigr)(X^2).
\]
所以乘法后的偶分支不是 \(AC\)，而是 \(AC+YBD\)；奇分支也不是 \(BD\)，而是 \(AD+BC\)。这说明 coefficient split 给出的是一个二维 module decomposition，而不是把大环乘法变成 leaf-wise multiplication 的 direct-product decomposition。
\end{example}

一般 \(k\) 的情况与上述二维例子相同：\(\Split_k\) 是一个线性 module decomposition，与加法兼容，但乘法会在不同 residue components 之间产生交叉项。因此，本文使用 coefficient split 的原因不是为了把所有同态运算都变成 leaf-wise independent computation，而是因为 CKKS bootstrapping core 具有特殊结构。更具体地说，bootstrapping core 可以写成“线性变换 \(\CtwoS\)、逐坐标非线性 \(\EvalMod\)、反向线性变换 \(\StwoC\)”的形式；coefficient split 与前后两个线性变换兼容，中间的 \(\EvalMod\) 又是逐坐标作用。这才是后文能够证明 leaf-wise bootstrapping factorization 的原因。

\subsection{子环嵌入与压缩的 slot 几何}
\label{subsec:subring-embedding-slot-geometry}

在讨论一般 coefficient split 的 slot 语义之前，先看一个更简单的特殊情况：消息本来就在子环嵌入像中。设 \(N=2n\)，小环消息为 \(f(Y)\in\mathcal R_n\)，其嵌入到大环后的消息为 \(f(X^2)\in\mathcal R_{2n}\)。若小环 slot vector 为 \(z=(f(\theta_0),\ldots,f(\theta_{n-1}))\)，那么大环中 \(f(X^2)\) 的 slot values 是 \(f(\alpha_i^2)\)。

在成对排列下，\(\alpha_{j+n}=-\alpha_j\)，且 \(\alpha_j^2=\alpha_{j+n}^2=\theta_j\)。因此
\[
f(\alpha_j^2)=f(\theta_j),\qquad
f(\alpha_{j+n}^2)=f(\theta_j).
\]
所以大环 slot vector 具有重复结构：
\[
\operatorname{Slots}_{2n}(f(X^2))=(z,z),
\]
这里等式忽略了实现中的固定 layout permutation。相反，从 \(f(X^2)\) 压缩回 \(f(Y)\) 在 slot 层面就是
\[
(z,z)\longmapsto z.
\]

\begin{example}[子环嵌入导致重复 slots]
取 \(n=2,N=4\)，令 \(f(Y)=b_0+b_1Y\)。设小环 roots 为 \(\theta_0,\theta_1\)，大环 roots 成对为 \(\alpha_0,-\alpha_0\) 和 \(\alpha_1,-\alpha_1\)，其中 \(\alpha_j^2=\theta_j\)。则
\[
f(X^2)(\alpha_j)=f(\theta_j),\qquad
f(X^2)(-\alpha_j)=f(\theta_j).
\]
因此大环中每一对 paired slots 完全相同。这说明 \(f(Y)\mapsto f(X^2)\) 不是把信息变复杂，而是在大环中重复存放同一份小环 slot 信息。
\end{example}

这个特殊情况对理解 ring compression 很有帮助。如果输入密文满足 \(\ct_{2n}\vDash_{\Delta,2n} f(X^2)\)，那么理想的压缩输出应满足 \(\ct_n\vDash_{\Delta,n} f(Y)\)。在 slot 语义上，它把重复结构 \((z,z)\) 压回一份 \(z\)。如果输入噪声写成 \(e(X)=e_0(X^2)+Xe_1(X^2)\)，那么压缩后保留的主要是偶部噪声 \(e_0\)，并额外加入 key-switching noise。这可以直观理解为：压缩保留 paired slots 的对称部分，丢弃反对称部分。

\subsection{Coefficient split 在 slot 视角下发生什么}
\label{subsec:coeff-split-slot-view}

现在回到一般消息 \(F(X)\in\mathcal R_{2n}\)。写成
\[
F(X)=A(X^2)+XB(X^2).
\]
设 \(\theta_j=\alpha_j^2\)，并记 \(a_j=A(\theta_j)\)、\(b_j=B(\theta_j)\)。则成对的大环 slots 满足
\[
F(\alpha_j)=a_j+\alpha_j b_j,\qquad
F(-\alpha_j)=a_j-\alpha_j b_j.
\]
因此，coefficient split 输出的 \(A(Y)\) 和 \(B(Y)\) 在小环 slots 中对应
\[
a_j=\frac{F(\alpha_j)+F(-\alpha_j)}{2},\qquad
b_j=\frac{F(\alpha_j)-F(-\alpha_j)}{2\alpha_j}.
\]
这说明 coefficient split 在 slot 视角下不是选择原始大环 slots，而是在每一对 paired slots 上做局部 butterfly：它提取对称部分 \(a_j\)，以及归一化后的反对称部分 \(b_j\)。

\begin{proposition}[coefficient split 不是 slot selection]
对于 \(N=2n\)，coefficient split \(F(X)=A(X^2)+XB(X^2)\) 在 slot 语义下将一对 big-ring slots \((F(\alpha_j),F(-\alpha_j))\) 变为 \((A(\theta_j),B(\theta_j))\)。该变换为局部线性变换，而不是坐标选择。
\end{proposition}

\begin{proof}
由 \(F(\alpha_j)=A(\theta_j)+\alpha_jB(\theta_j)\) 与 \(F(-\alpha_j)=A(\theta_j)-\alpha_jB(\theta_j)\)，相加得到 \(A(\theta_j)=(F(\alpha_j)+F(-\alpha_j))/2\)，相减得到 \(B(\theta_j)=(F(\alpha_j)-F(-\alpha_j))/(2\alpha_j)\)。
\end{proof}

一般 \(k\) 的情况可以按 fiber 理解。固定一个小环 root \(\theta\)，取其上方的 \(k\) 个大环 roots \(\alpha_{\theta,0},\ldots,\alpha_{\theta,k-1}\)，满足 \(\alpha_{\theta,\ell}^k=\theta\)。由 \(P(X)=\sum_{r=0}^{k-1}X^rP_r(X^k)\)，对每个 \(\ell\) 有
\[
P(\alpha_{\theta,\ell})=\sum_{r=0}^{k-1}\alpha_{\theta,\ell}^{\,r}P_r(\theta).
\]
记 \(z_\theta=(P(\alpha_{\theta,0}),\ldots,P(\alpha_{\theta,k-1}))^T\)，\(u_\theta=(P_0(\theta),\ldots,P_{k-1}(\theta))^T\)。则
\[
z_\theta=V_\theta u_\theta,
\qquad
(V_\theta)_{\ell,r}=\alpha_{\theta,\ell}^{\,r}.
\]
因此 \(u_\theta=V_\theta^{-1}z_\theta\)。这说明一般 \(k\) 下，coefficient split 在 slot 视角中对应每个 fiber 上的局部 Vandermonde 逆变换。

\begin{example}[从 paired slots 到偶/奇分支]
仍取 \(N=2n\)。如果某一对大环 slots 为 \((u+\alpha v,u-\alpha v)\)，那么 coefficient split 后的两个小环分支在对应小环 slot 上就是 \(u\) 和 \(v\)。如果输入本来是子环嵌入 \(f(X^2)\)，则 \(v=0\)，paired slots 退化为 \((u,u)\)。这正是上一小节中重复结构 \((z,z)\) 的特殊情形。
\end{example}

\subsection{Trace-like projection 的补充视角}
\label{subsec:trace-like-projection}

上面的 even/odd decomposition 也可以用 trace-like projection 来理解。对于 \(N=2n\)，固定小环 \(\mathbb C(X^2)\) 的非平凡自同构可以理解为 \(X\mapsto -X\)。因此
\[
\frac{F(X)+F(-X)}{2}=A(X^2),
\qquad
\frac{F(X)-F(-X)}{2X}=B(X^2).
\]
第一式提取偶部，第二式提取归一化后的奇部。

\begin{example}[trace 提取偶部]
设 \(F(X)=a_0+a_1X+a_2X^2+a_3X^3\)。则
\[
\frac{F(X)+F(-X)}{2}=a_0+a_2X^2=A(X^2),
\]
而
\[
\frac{F(X)-F(-X)}{2X}=a_1+a_3X^2=B(X^2).
\]
这说明 coefficient split 的 \(A,B\) 可以看成由对 \(X\mapsto -X\) 的对称/反对称投影得到。
\end{example}

该视角有助于理解“为什么偶部落在小环里”。不过本文的主算法并不是直接把 trace 作为密文操作，而是通过 HERMES-style ring switching 获得 leaf ciphertexts。因此，本小节只作为语义解释。

\subsection{Slot split 与 slot-preserving split}
\label{subsec:slot-split-preserving}

与 coefficient split 相对，slot split 的目标是直接控制 leaf ciphertexts 的 slot 语义。给定 \(\ct_N\vDash_{\Delta,N}P(X)\)，其 big-ring slot vector 可写为 \(z=F_Np\)，其中 \(p\) 是 \(P\) 的系数向量，\(F_N\) 是大环 evaluation matrix。如果我们希望某个小环密文 \(\ct_q\) 表示 \(Q(Y)\)，并且 \(Q\) 的小环 slot vector 是指定的 \(y=Tz\)，那么 \(Q\) 的系数向量 \(q\) 必须满足 \(F_nq=y\)。因此
\[
q=F_n^{-1}TF_Np.
\]
这就是一般 slot split 的系数形式。这里 \(T\) 可以是 slot selection，也可以是任意 slot-level linear transform。

这个公式说明：slot split 不是简单的 coefficient residue-class selection。除非 \(T\) 与 \(F_N,F_n\) 有特殊结构，否则 \(F_n^{-1}TF_N\) 通常是稠密线性变换。在密文层面，它意味着需要执行额外的 homomorphic linear transform。因此，slot split 更通用，但通常更贵。

Slot-preserving split 是 slot split 的一个重要特例。仍以 \(N=2n\) 为例。我们希望构造两个小环多项式 \(Q^+(Y)\) 和 \(Q^-(Y)\)，使得
\[
Q^+(\theta_j)=F(\alpha_j),\qquad
Q^-(\theta_j)=F(-\alpha_j).
\]
而 coefficient split 只给出 \(A(Y)\) 和 \(B(Y)\)。由 \(F(\alpha_j)=A(\theta_j)+\alpha_jB(\theta_j)\) 可知，如果存在小环多项式 \(h(Y)\) 满足 \(h(\theta_j)=\alpha_j\)，则可定义
\[
Q^+(Y)=A(Y)+h(Y)B(Y),\qquad
Q^-(Y)=A(Y)-h(Y)B(Y).
\]
于是 \(Q^+(\theta_j)=F(\alpha_j)\)，且 \(Q^-(\theta_j)=F(-\alpha_j)\)。

在密文层面，如果 \(\ct_A\vDash_{\Delta,n}A(Y)\) 和 \(\ct_B\vDash_{\Delta,n}B(Y)\)，则 slot-preserving split 的第二步需要计算
\[
\ct^+=\ct_A+\operatorname{MulPt}_{h(Y)}(\ct_B),
\qquad
\ct^-=\ct_A-\operatorname{MulPt}_{h(Y)}(\ct_B).
\]
这一步不是 coefficient split 本身的一部分，而是额外的 leaf-side homomorphic linear transform。

用矩阵形式表示，令 \(E\) 与 \(O\) 分别为偶、奇系数选择矩阵，令 \(D=\operatorname{diag}(\alpha_0,\ldots,\alpha_{n-1})\)，并令 \(H=F_n^{-1}DF_n\)。则 coefficient split 是 \(p\mapsto(Ep,Op)\)，而 pairwise slot-preserving split 是
\[
p\mapsto(Ep+HOp,\;Ep-HOp).
\]
因此，slot-preserving split 比 coefficient split 多了一个 \(H\)，也就是小环上乘以平方根选择多项式 \(h(Y)\) 的线性变换。

\begin{example}[slot-preserving split 的额外代价]
如果 \(F(X)=A(X^2)+XB(X^2)\)，coefficient split 只输出表示 \(A\) 和 \(B\) 的两个 leaf ciphertexts。若希望两个 leaf ciphertexts 的 slots 分别等于 \(F(\alpha_j)\) 和 \(F(-\alpha_j)\)，则还必须从 \(A,B\) 构造 \(A+hB\) 与 \(A-hB\)。这一步在明文语义上很简单，但在密文层面需要额外的明文乘法和加减法；一般 slot selection \(q=F_n^{-1}TF_Np\) 则可能对应更稠密的同态线性变换。
\end{example}

\subsection{Slot-level split as representation switching}
\label{subsec:slot-level-representation-switching}

还有一种常见的 slot-level split 来自 CKKS 的表示切换。它不是按 \(Y=X^k\) 做 ring switching，而是在同一个 CKKS 编码语义下，把一个 slot 向量中携带的系数信息同态地分成两个密文。

设
\[
p(Y)=\sum_{i=0}^{2n-1}p_iY^i,\qquad
\mathbf p_0=(p_0,\ldots,p_{n-1}),\qquad
\mathbf p_1=(p_n,\ldots,p_{2n-1}).
\]
定义复向量 \(\mathbf w=\mathbf p_0+\mathbf i\,\mathbf p_1\)，并令 \(\mathbf z=U_n\mathbf w\)。在 CKKS 编码语义下，一个密文可以表示 \(\operatorname{Ecd}(\mathbf z)\)，也就是表示由 \(p(Y)\) 编码得到的消息。若对该密文同态乘以 \(U_n^{-1}\)，则 slot 语义变为 \(\mathbf p_0+\mathbf i\,\mathbf p_1\)。随后通过共轭可分离实部和虚部：
\[
\frac12(\mathbf w+\overline{\mathbf w})=\mathbf p_0,
\qquad
\frac{1}{2\mathbf i}(\mathbf w-\overline{\mathbf w})=\mathbf p_1.
\]
密文层面相应地得到两个密文，分别表示 \(\operatorname{Ecd}(\mathbf p_0)\) 和 \(\operatorname{Ecd}(\mathbf p_1)\)。

反向的 SlotToCoeff 则从两个密文开始：一个表示 \(\operatorname{Ecd}(\mathbf p_0)\)，另一个表示 \(\operatorname{Ecd}(\mathbf p_1)\)。先同态构造 \(\mathbf p_0+\mathbf i\,\mathbf p_1\)，再乘以 \(U_n\)，得到表示 \(\operatorname{Ecd}(\mathbf z)\) 的密文。这个过程可以理解为把“slot 中装着系数信息”的表示切换回“一个多项式被编码”的表示。

\begin{example}[一个密文拆成两个 slot 密文]
假设输入密文表示 \(\operatorname{Ecd}(\mathbf z)\)，其中 \(\mathbf z=U_n(\mathbf p_0+\mathbf i\,\mathbf p_1)\)。同态乘 \(U_n^{-1}\) 后，密文表示 \(\operatorname{Ecd}(\mathbf p_0+\mathbf i\,\mathbf p_1)\)。再做
\[
\ct_0=\frac12(\ct'+\operatorname{Conj}(\ct')),
\qquad
\ct_1=\frac{1}{2\mathbf i}(\ct'-\operatorname{Conj}(\ct')),
\]
即可分别得到表示 \(\mathbf p_0\) 和 \(\mathbf p_1\) 的两个密文。这个过程是 slot-level representation switching，而不是本文主算法中的 HERMES-style ring switching。
\end{example}

这类 slot-level split 对解释 CKKS 的 CoeffToSlot/SlotToCoeff 很有用，但它与本文的 coefficient/module split 不同。本文的目标是把一个大环 ciphertext 转换为多个小环 ciphertexts，使 leaf ciphertexts 分别表示 \(P_r(Y)\)。而 slot-level split 通常仍是在编码语义中做同态线性变换，并不自动改变 ciphertext 所在的环维度。

\subsection{Two-step splitting interpretation}
\label{subsec:two-step-splitting}

上述几种 split 可以统一为一个 two-step view。

第一步是 ciphertext-level coefficient/module decomposition。输入是大环密文 \(\ct_N\vDash_{\Delta,N}P(X)\)。通过 ring switching，输出 \(k\) 个 leaf ciphertexts \(\ct_{n,r}\)，并使得 \(\ct_{n,r}\vDash_{\Delta,n}P_r(Y)\)。这一步改变了密文所在的环和密钥语义，是本文加速的基础。

第二步是可选的 slot-coordinate transformation。如果下游任务要求 leaf ciphertexts 的 slots 具有特定的 big-ring slot 语义，则需要在第一步之后继续执行额外同态线性变换。对于 \(k=2\)，从 coefficient split 到 pairwise slot-preserving split 的第二步是 \(A,B\mapsto A\pm hB\)。对于一般 slot selection，它由 \(q=F_n^{-1}TF_Np\) 描述。对于 CKKS 表示切换场景，它也可能表现为 \(U_n^{-1}\)、conjugation、\(U_n\) 这一类 slot-level linear transforms。

本文的 bootstrapping 分解只使用第一步，不执行第二步。原因是 CKKS bootstrapping core 并不要求 split 后立即恢复原始 big-ring slots。它需要的是：经过 \(\CtwoS\) 后，slots 中装入 coefficient-like values，然后 \(\EvalMod\) 对这些 values 逐坐标作用。Coefficient split 后直接接 leaf \(\CtwoS_n\)，恰好得到大环 \(\CtwoS_N\) 输出的 coefficient residue-class 子序列。因此，slot-preserving transformation 不是本文主算法的一部分。

\subsection{为什么 bootstrapping 分解选择 coefficient split}
\label{subsec:why-coeff-split-for-bts}

现在可以解释本文为什么选择 coefficient split，而不是 slot-preserving split。

CKKS bootstrapping core 的逻辑不是“在 split 后保留原始 big-ring slots 并继续计算”。传统 core 的结构是先通过 \(\CtwoS\) 把待模约减的 coefficient-like values 放入 slots，然后用 \(\EvalMod\) 逐坐标去掉 \(q_0I_i\)，最后通过 \(\StwoC\) 回到系数表示。因此，分解后真正需要对齐的是 \(\CtwoS\) 之后的 coefficient-packed coordinates，而不是 split 之后的 canonical slots。

Coefficient split 虽然在 canonical slot 视角下会混合 big-ring slots，但它有一个关键优点：若先对密文执行 ring-switch-down 得到 \(\ct_{n,r}\vDash_{\Delta,n}P_r(Y)\)，再对每个 leaf 执行 \(\CtwoS_n\)，则 leaf 的 \(\CtwoS_n\) 输出正是 \(P_r\) 的系数，也就是大环 \(P\) 的 residue-class 系数子序列。换句话说，\(\Split_k+\bigoplus_r\CtwoS_n\) 在语义上等价于先做 top \(\CtwoS_N\)，再把输出按 coefficient residue classes 分组。

这解释了为什么本文不需要先做 slot-preserving split。Slot-preserving split 试图恢复 \(P(\alpha)\)、\(P(-\alpha)\) 这类 big-ring slot values；但 bootstrapping 的 \(\EvalMod\) 阶段需要处理的是 \(\CtwoS\) 后的 coefficient-packed values。额外恢复 big-ring slots 不仅不是必要步骤，还会引入额外线性变换成本。

\subsection{本节小结}
\label{subsec:coeff-slot-summary}

本节的结论可以概括为四点。第一，本文讨论的 split 是密文语义下的 split：给定 \(\ct_N\vDash_{\Delta,N}P(X)\)，目标是通过 ring switching 得到 \(\ct_{n,r}\vDash_{\Delta,n}P_r(Y)\)，而不是直接读取 \(P\) 或其系数。第二，coefficient/module split 在系数层面是简单的 residue-class decomposition，但在 slot 视角下不是 slot selection；对于 \(k=2\)，它提取 paired slots 的平均部分和加权差值。第三，如果希望拆分后保留原始 big-ring slot 语义，就需要 slot split 或 slot-preserving split，这通常需要额外同态线性变换，例如 \(A\pm hB\) 或更一般的 \(F_n^{-1}TF_N\)。第四，本文的 bootstrapping factorization 选择 coefficient split，是因为 \(\Split_k\) 后接 leaf \(\CtwoS_n\) 正好获得 top \(\CtwoS_N\) 输出的 coefficient residue-class 子序列。下一节将把这一点形式化为 Split--C2S compatibility，并进一步证明 \(\EvalMod\) 与 \(\StwoC+\Merge\) 的兼容性。

\subsection{本节小结}
\label{subsec:coeff-slot-summary}

本节的结论可以总结如下。

第一，本文讨论的 split 是密文语义下的 split。给定 \(\ct_N\rep_N P(X)\)，目标是通过 ring switching 得到 \(\ct_{n,r}\rep_n P_r(Y)\)，而不是直接读取 \(P\) 或其系数。

第二，coefficient/module split 在系数层面是简单的 residue-class decomposition：\(p\mapsto(S_0p,\ldots,S_{k-1}p)\)。但在 slot 视角下，它不是 slot selection；对于 \(k=2\)，它把 \((P(\alpha),P(-\alpha))\) 变成 \((A(\theta),B(\theta))\)，即平均和加权差值；对于一般 \(k\)，它对应每个 fiber 上的局部 Vandermonde 逆变换。

第三，如果希望拆分后保留原始 big-ring slot 语义，就需要 slot split 或 slot-preserving split。这通常需要额外的同态线性变换，例如 \(A\pm hB\)，或更一般的 \(F_n^{-1}TF_N\)。

第四，本文的 bootstrapping factorization 选择 coefficient split，是因为 \(\Split_k\) 后接 leaf \(\CtwoS_n\) 正好获得 top \(\CtwoS_N\) 输出的 coefficient residue-class 子序列。下一节将把这一点形式化为 Split--C2S compatibility，并进一步证明 \(\EvalMod\) 与 \(\StwoC+\Merge\) 的兼容性。
\section{主要理论结果}
\label{sec:main-theoretical-results}

上一节说明了本文使用的 split 是密文语义下的 coefficient/module split：若输入密文满足 \(\ct_N\vDash_{\Delta,N}P(X)\)，则理想的 ring-switch-down 输出应满足 \(\ct_{n,r}\vDash_{\Delta,n}P_r(Y)\)，其中 \(P(X)=\sum_{r=0}^{k-1}X^rP_r(X^k)\)。本节在消息语义和 normalized logical layout 下证明本文的核心代数结论：传统 CKKS bootstrapping core 可以沿该 coefficient split 分解为多个 leaf bootstrapping core。

为避免实现细节干扰，本节暂时忽略 Lattigo 中的 bit-reversal、blocked layout、real/imag branches、half-slot convention 和 scaling convention。所有这些实现层细节都可以视作固定 layout permutation 或 implementation scaling，并在实现章节单独讨论。本节只证明理想语义下的结构性等式。

\subsection{Notations}
\label{subsec:theory-notations}

固定 \(N=kn\)。令 \(\mathsf{Can}_d\) 表示环维度为 \(d\) 时的 canonical-slot 坐标空间，令 \(\mathsf{Coeff}_d\) 表示对应的 coefficient-packed 坐标空间。若 \(P\in\mathcal R_d\)，记 \(\operatorname{can}_d(P)\in\mathsf{Can}_d\) 为 \(P\) 的 canonical-slot 向量，记 \(\operatorname{coeff}_d(P)\in\mathsf{Coeff}_d\) 为 \(P\) 的系数向量。

本节把 \(\CtwoS_d:\mathsf{Can}_d\to\mathsf{Coeff}_d\) 视为理想的 coefficient-to-slot 反向表示切换：它把 canonical-slot 坐标中的消息转为 EvalMod 所作用的 coefficient-packed 坐标。相应地，\(\StwoC_d:\mathsf{Coeff}_d\to\mathsf{Can}_d\) 是反向线性变换。在 full canonical embedding 的 normalized model 下，可以把 \(\CtwoS_d\) 理解为 evaluation matrix 的逆，把 \(\StwoC_d\) 理解为 evaluation matrix；但是本文的证明只使用二者的语义关系：
\[
\CtwoS_d(\operatorname{can}_d(P))=\operatorname{coeff}_d(P),
\qquad
\StwoC_d(\operatorname{coeff}_d(P))=\operatorname{can}_d(P).
\]
这里单独展示该式是因为它是后续所有交换图的语义基础。

对任意大环系数向量 \(p=(a_0,\ldots,a_{N-1})^T\in\mathsf{Coeff}_N\)，定义 residue-class 选择算子 \(S_r:\mathsf{Coeff}_N\to\mathsf{Coeff}_n\) 为
\[
S_rp=(a_r,a_{r+k},a_{r+2k},\ldots,a_{r+(n-1)k})^T,
\qquad r=0,\ldots,k-1.
\]
令 \(S_kp=(S_0p,\ldots,S_{k-1}p)\)。在 normalized layout 下，\(S_k\) 就是把长度为 \(N\) 的 coefficient-packed vector 按 residue class 分组成 \(k\) 个长度为 \(n\) 的 leaf vectors。

另一方面，上一节定义的明文 split 为
\[
P(X)=\sum_{r=0}^{k-1}X^rP_r(X^k),
\qquad
\Split_k(P)=(P_0,\ldots,P_{k-1}).
\]
在 canonical-slot 坐标中，我们也记 \(\Split_k\) 为由该消息 split 诱导的语义 map：
\[
\Split_k(\operatorname{can}_N(P))
=
(\operatorname{can}_n(P_0),\ldots,\operatorname{can}_n(P_{k-1})).
\]
类似地，\(\Merge_k\) 在 canonical-slot 语义下表示先将 leaf 消息 \(P_0,\ldots,P_{k-1}\) 组装成 \(\sum_rX^rP_r(X^k)\)，再取大环 canonical-slot 表示。

最后，用 \(\Pi_k:\mathsf{Coeff}_N\to\bigoplus_{r=0}^{k-1}\mathsf{Coeff}_n\) 表示 top coefficient-packed vector 到 leaf coefficient-packed vectors 的 layout map。在 normalized logical layout 下，\(\Pi_k=S_k\)。在具体实现中，\(\Pi_k\) 还可能包含 bit-reversal、block ordering 或 real/imag branch ordering；这些都不影响本节的代数结论。

\begin{example}[符号在 \(k=2\) 时的含义]
若 \(N=2n\)，任意 \(P(X)\) 可写成 \(P(X)=A(X^2)+XB(X^2)\)。若 \(p=(a_0,a_1,a_2,a_3,\ldots)^T\)，则 \(S_0p=(a_0,a_2,a_4,\ldots)^T\) 是 \(A(Y)\) 的系数向量，\(S_1p=(a_1,a_3,a_5,\ldots)^T\) 是 \(B(Y)\) 的系数向量。因此 \(\Pi_2p=(S_0p,S_1p)\) 就是把 top coefficient-packed vector 分成偶系数分支和奇系数分支。
\end{example}

\subsection{引理与定理陈述}
\label{subsec:lemmas-theorems-statement}

本节的核心结论由三个兼容性引理组成。它们分别对应 CKKS bootstrapping core 的三段：\(\CtwoS\)、\(\EvalMod\)、\(\StwoC\)。直观上，第一条引理说明 split 后做 leaf \(\CtwoS_n\) 等价于先做 top \(\CtwoS_N\) 再分组；第二条引理说明逐坐标的 \(\EvalMod\) 可以穿过这种分组；第三条引理说明 leaf \(\StwoC_n\) 后再 merge 等价于 top \(\StwoC_N\)。

\begin{lemma}[Split--C2S compatibility]
\label{lem:split-c2s}
对于 \(N=kn\)，在上述 normalized layout convention 下，
\(\left(\bigoplus_{r=0}^{k-1}\CtwoS_n\right)\circ\Split_k=\Pi_k\circ\CtwoS_N.\)
\end{lemma}

\begin{lemma}[EvalMod compatibility]
\label{lem:evalmod}
若 \(\EvalMod_N=f^{\oplus N}\)、\(\EvalMod_n=f^{\oplus n}\)，且二者由同一个 coordinate-wise scalar function \(f\) 定义，则
\(\left(\bigoplus_{r=0}^{k-1}\EvalMod_n\right)\circ\Pi_k=\Pi_k\circ\EvalMod_N.\)
\end{lemma}

\begin{lemma}[S2C--Merge compatibility]
\label{lem:s2c-merge}
在相同 layout convention 下，
\(\Merge_k\circ\left(\bigoplus_{r=0}^{k-1}\StwoC_n\right)\circ\Pi_k=\StwoC_N.\)
\end{lemma}

\begin{theorem}[CKKS bootstrapping core factorization]
\label{thm:ckks-core-factorization}
在上述 normalized layout convention 下，CKKS bootstrapping core 满足
\begin{equation}
\label{eq:ckks-core-factorization}
\StwoC_N\circ\EvalMod_N\circ\CtwoS_N
=
\Merge_k\circ
\left(
\bigoplus_{r=0}^{k-1}
\StwoC_n\circ\EvalMod_n\circ\CtwoS_n
\right)
\circ\Split_k.
\end{equation}
\end{theorem}

式 \eqref{eq:ckks-core-factorization} 的意义是：理想模型中，一个 top-level bootstrapping core 可以被替换为 split、\(k\) 个 leaf core、再 merge 的组合。注意，该等式不是说 split 后 leaf canonical slots 保留了原始 big-ring slots；上一节已经说明 coefficient split 在 slot 视角下通常是局部 Vandermonde 型变换。这里真正起作用的是：split 后接 leaf \(\CtwoS_n\) 会进入 coefficient-packed coordinates，而这些 coordinates 正是 top \(\CtwoS_N\) 输出的 residue-class 子序列。

\subsection{Split--C2S compatibility 的证明}
\label{subsec:proof-split-c2s}

\begin{proof}
任取一个大环消息多项式 \(P(X)=\sum_{i=0}^{N-1}a_iX^i\)，记其大环 canonical-slot 向量为 \(z=\operatorname{can}_N(P)\)，系数向量为 \(p=\operatorname{coeff}_N(P)=(a_0,\ldots,a_{N-1})^T\)。根据 \(\CtwoS_N\) 的语义定义，有 \(\CtwoS_N(z)=p\)。

现在写出 \(P\) 的 coefficient/module split：
\[
P(X)=\sum_{r=0}^{k-1}X^rP_r(X^k),
\qquad
P_r(Y)=\sum_{j=0}^{n-1}a_{r+jk}Y^j.
\]
因此，第 \(r\) 个 leaf 消息 \(P_r\) 的系数向量正是 \(S_rp\)。

先看左侧路径。输入 \(z=\operatorname{can}_N(P)\) 经过 \(\Split_k\) 后，按照定义得到
\[
\Split_k(z)
=
(\operatorname{can}_n(P_0),\ldots,\operatorname{can}_n(P_{k-1})).
\]
再对每个 leaf 应用 \(\CtwoS_n\)，第 \(r\) 个分支变为
\[
\CtwoS_n(\operatorname{can}_n(P_r))
=
\operatorname{coeff}_n(P_r)
=
S_rp.
\]
因此左侧整体输出为
\[
\begin{aligned}
\left(
\left(\bigoplus_{r=0}^{k-1}\CtwoS_n\right)
\circ\Split_k
\right)(z)
&=
(\CtwoS_n(\operatorname{can}_n(P_0)),\ldots,\CtwoS_n(\operatorname{can}_n(P_{k-1})))\\
&=
(\operatorname{coeff}_n(P_0),\ldots,\operatorname{coeff}_n(P_{k-1}))\\
&=
(S_0p,\ldots,S_{k-1}p)\\
&=
S_kp.
\end{aligned}
\]

再看右侧路径。输入 \(z\) 先经过 top-level \(\CtwoS_N\)，得到 \(p\)。随后 \(\Pi_k\) 将 \(p\) 按 residue class 分组。在 normalized layout 下，\(\Pi_kp=S_kp\)。所以右侧输出为
\[
\begin{aligned}
(\Pi_k\circ\CtwoS_N)(z)
&=
\Pi_k(\CtwoS_N(z))\\
&=
\Pi_kp\\
&=
S_kp.
\end{aligned}
\]

左右两条路径对任意 \(P\) 的 canonical-slot 输入 \(z\) 都得到同一个 grouped coefficient vector \(S_kp\)。因此
\[
\left(\bigoplus_{r=0}^{k-1}\CtwoS_n\right)\circ\Split_k
=
\Pi_k\circ\CtwoS_N.
\]
\end{proof}

\begin{example}[\(k=2\) 时的 Split--C2S compatibility]
令 \(P(X)=A(X^2)+XB(X^2)\)。左侧路径先 split 得到 \(A(Y)\) 和 \(B(Y)\)，再分别做 leaf \(\CtwoS_n\)，得到 \(A\) 和 \(B\) 的系数向量，即 \((a_0,a_2,\ldots)\) 与 \((a_1,a_3,\ldots)\)。右侧路径先做 top \(\CtwoS_N\)，得到 \(P\) 的完整系数向量 \((a_0,a_1,a_2,a_3,\ldots)\)，再按偶/奇下标分组。两条路径输出相同。这一例子也说明：引理不是说 split 保留了原始 big slots，而是说 split 后接 \(\CtwoS_n\) 恰好得到 top \(\CtwoS_N\) 后的系数子序列。
\end{example}

\subsection{EvalMod compatibility 的证明}
\label{subsec:proof-evalmod}

\begin{proof}
令 \(x=(x_0,\ldots,x_{N-1})\in\mathsf{Coeff}_N\) 为 top \(\CtwoS_N\) 后的 coefficient-packed vector。根据 \(\Pi_k\) 的定义，\(\Pi_kx=(S_0x,\ldots,S_{k-1}x)\)，其中第 \(r\) 个 leaf 的第 \(j\) 个坐标为 \((S_rx)_j=x_{r+jk}\)。

左侧路径先分组，再在每个 leaf 上做 \(\EvalMod_n=f^{\oplus n}\)。因此左侧输出在第 \(r\) 个 leaf、第 \(j\) 个坐标处为
\[
\begin{aligned}
\left[
\left(\bigoplus_{r=0}^{k-1}\EvalMod_n\right)(\Pi_kx)
\right]_{r,j}
&=
\left[\EvalMod_n(S_rx)\right]_j\\
&=
f((S_rx)_j)\\
&=
f(x_{r+jk}).
\end{aligned}
\]

右侧路径先在 top vector 上做 \(\EvalMod_N=f^{\oplus N}\)，得到
\[
\EvalMod_N(x)=(f(x_0),f(x_1),\ldots,f(x_{N-1})).
\]
再由 \(\Pi_k\) 按 residue class 分组。因此右侧输出在第 \(r\) 个 leaf、第 \(j\) 个坐标处为
\[
\begin{aligned}
\left[
\Pi_k(\EvalMod_N(x))
\right]_{r,j}
&=
\left[\EvalMod_N(x)\right]_{r+jk}\\
&=
f(x_{r+jk}).
\end{aligned}
\]

左右两侧每个坐标都相等，因此
\[
\left(\bigoplus_{r=0}^{k-1}\EvalMod_n\right)\circ\Pi_k
=
\Pi_k\circ\EvalMod_N.
\]
\end{proof}

这个证明没有使用 \(f\) 的线性性，只使用了同一个 scalar function \(f\) 在每个坐标上独立作用。因此即使 \(\EvalMod\) 是非线性的，只要它是 coordinate-wise map，它仍然可以穿过这种 residue-class grouping。这一点正是本文能够把非线性 bootstrapping core 分解到 leaf 上的关键。

\begin{example}[逐坐标函数为什么可以穿过分组]
设 \(N=4,k=2,n=2\)，并令 \(x=(x_0,x_1,x_2,x_3)\)。分组后得到 \((x_0,x_2)\) 和 \((x_1,x_3)\)。先分组再应用 \(f\)，得到 \((f(x_0),f(x_2))\) 和 \((f(x_1),f(x_3))\)。先对 \(x\) 整体逐坐标应用 \(f\)，再分组，得到的也是同一结果。因此这里不需要 \(f(a+b)=f(a)+f(b)\) 之类的性质。
\end{example}

\subsection{S2C--Merge compatibility 的证明}
\label{subsec:proof-s2c-merge}

\begin{proof}
任取 EvalMod 后的大环 coefficient-packed vector \(y=(y_0,\ldots,y_{N-1})\in\mathsf{Coeff}_N\)。经过 \(\Pi_k\) 分组后，第 \(r\) 个 leaf 得到
\[
S_ry=(y_r,y_{r+k},\ldots,y_{r+(n-1)k}).
\]
将 \(S_ry\) 解释为小环多项式
\[
Y_r(Y)=\sum_{j=0}^{n-1}y_{r+jk}Y^j
\]
的系数向量。

先看左侧路径。对每个 leaf 应用 \(\StwoC_n\) 后，第 \(r\) 个 leaf 得到 \(\operatorname{can}_n(Y_r)\)。随后 \(\Merge_k\) 按 coefficient/module assembly 的规则把这些 leaf 消息合成一个大环消息：
\[
Y(X)=\sum_{r=0}^{k-1}X^rY_r(X^k).
\]
把 \(Y_r\) 展开，得到
\[
\begin{aligned}
Y(X)
&=
\sum_{r=0}^{k-1}X^r
\left(
\sum_{j=0}^{n-1}y_{r+jk}(X^k)^j
\right)\\
&=
\sum_{r=0}^{k-1}
\sum_{j=0}^{n-1}
y_{r+jk}X^{r+jk}.
\end{aligned}
\]
由于每个下标 \(i\in\{0,\ldots,N-1\}\) 都唯一写成 \(i=r+jk\)，上式正是
\[
Y(X)=\sum_{i=0}^{N-1}y_iX^i.
\]
因此，左侧路径最终输出的是系数向量为 \(y\) 的大环多项式的 canonical-slot 表示，即 \(\operatorname{can}_N(Y)\)。

右侧路径 \(\StwoC_N\) 则直接把 coefficient-packed vector \(y\) 解释为大环多项式 \(\sum_i y_iX^i\) 的系数向量，并输出其 canonical-slot 表示。因此右侧输出同样是 \(\operatorname{can}_N(Y)\)。所以
\[
\Merge_k\circ
\left(\bigoplus_{r=0}^{k-1}\StwoC_n\right)
\circ\Pi_k
=
\StwoC_N.
\]
\end{proof}

\begin{example}[\(k=2\) 时的 S2C--Merge compatibility]
设 \(y=(y_0,y_1,y_2,y_3)\)。分组得到偶分支 \((y_0,y_2)\) 和奇分支 \((y_1,y_3)\)，对应小环多项式 \(Y_0(Y)=y_0+y_2Y\) 和 \(Y_1(Y)=y_1+y_3Y\)。leaf \(\StwoC_n\) 后再 merge，得到
\[
Y_0(X^2)+XY_1(X^2)
=
y_0+y_1X+y_2X^2+y_3X^3,
\]
这正是 top \(\StwoC_N\) 会解释的同一个 coefficient-packed vector。
\end{example}

\subsection{CKKS core factorization 的证明}
\label{subsec:proof-ckks-core}

\begin{proof}
记 leaf core 为 \(\BTS_{n}^{\mathrm{core}}=\StwoC_n\circ\EvalMod_n\circ\CtwoS_n\)。从式 \eqref{eq:ckks-core-factorization} 的右侧出发，逐步使用前面三个兼容性引理：
\[
\begin{aligned}
&\Merge_k\circ
\left(
\bigoplus_{r=0}^{k-1}
\StwoC_n\circ\EvalMod_n\circ\CtwoS_n
\right)
\circ\Split_k\\
&\quad=
\Merge_k\circ
\left(\bigoplus_r\StwoC_n\right)
\circ
\left(\bigoplus_r\EvalMod_n\right)
\circ
\left(\bigoplus_r\CtwoS_n\right)
\circ\Split_k\\
&\quad=
\Merge_k\circ
\left(\bigoplus_r\StwoC_n\right)
\circ
\left(\bigoplus_r\EvalMod_n\right)
\circ
\Pi_k\circ\CtwoS_N\\
&\quad=
\Merge_k\circ
\left(\bigoplus_r\StwoC_n\right)
\circ
\Pi_k\circ\EvalMod_N\circ\CtwoS_N\\
&\quad=
\StwoC_N\circ\EvalMod_N\circ\CtwoS_N.
\end{aligned}
\]
第一步只是展开 direct sum 中的 leaf core；第二步使用 Lemma~\ref{lem:split-c2s}；第三步使用 Lemma~\ref{lem:evalmod}；第四步使用 Lemma~\ref{lem:s2c-merge}。因此 CKKS bootstrapping core 的 factorization 成立。
\end{proof}

这个定理说明，理想模型中的分解路径不是在 split 后恢复原始大环 slots，而是直接利用 \(\Split_k+\)leaf \(\CtwoS_n\) 获得大环 \(\CtwoS_N\) 输出的 coefficient 子序列。由于 \(\EvalMod\) 是逐坐标函数，这些子序列可以在 leaf 上独立处理；最后 leaf \(\StwoC_n\) 与 \(\Merge_k\) 再把结果重组成 top \(\StwoC_N\) 的输出。

从密文语义看，若 ring switching 实现了 \(\ct_N\vDash_{\Delta,N}P(X)\mapsto(\ct_{n,r}\vDash_{\Delta,n}P_r(Y))_r\)，那么 leaf bootstrapping core 就是在这些 \(P_r\) 上分别执行同一个理想 core。理论等式保证了这些 leaf core 的消息语义可以 merge 回 top core 的消息语义。ciphertext-level 的 key-switching noise、rescale error 和 implementation scaling 不属于本节的代数等式，而是在算法与实现章节中纳入误差分析。

\subsection{General split-compatible sandwich theorem}
\label{subsec:general-sandwich-theorem}

上面的 CKKS core factorization 其实是一个更一般结构命题的实例。这个抽象定理有助于说明本文证明的真正原因：不是 CKKS 的某个实现细节偶然成立，而是“前向线性变换--逐坐标非线性--反向线性变换”这种 sandwich structure 与 split/merge 兼容。

\begin{theorem}[General split-compatible sandwich theorem]
\label{thm:general-sandwich}
设 \(F_N=L_N^{(2)}\circ\Phi_N\circ L_N^{(1)}\)，其中 \(L_N^{(1)}\) 和 \(L_N^{(2)}\) 是线性映射，\(\Phi_N=\phi^{\oplus N}\) 是逐坐标映射。设 leaf 电路为 \(F_n=L_n^{(2)}\circ\Phi_n\circ L_n^{(1)}\)，其中 \(\Phi_n=\phi^{\oplus n}\)。若存在 layout map \(\Pi_k\)，使得
\[
\left(\bigoplus_{r=0}^{k-1}L_n^{(1)}\right)\circ\Split_k
=
\Pi_k\circ L_N^{(1)}
\]
且
\[
\Merge_k\circ
\left(\bigoplus_{r=0}^{k-1}L_n^{(2)}\right)
\circ\Pi_k
=
L_N^{(2)},
\]
则
\[
F_N
=
\Merge_k\circ
\left(\bigoplus_{r=0}^{k-1}F_n\right)
\circ\Split_k.
\]
\end{theorem}

\begin{proof}
从右侧出发，先展开 leaf 电路：
\[
\begin{aligned}
&\Merge_k\circ
\left(\bigoplus_r F_n\right)
\circ\Split_k\\
&\quad=
\Merge_k\circ
\left(\bigoplus_r L_n^{(2)}\right)
\circ
\left(\bigoplus_r \Phi_n\right)
\circ
\left(\bigoplus_r L_n^{(1)}\right)
\circ\Split_k.
\end{aligned}
\]
由第一个兼容条件，\((\bigoplus_r L_n^{(1)})\circ\Split_k=\Pi_k\circ L_N^{(1)}\)。代入得
\[
\Merge_k\circ
\left(\bigoplus_r L_n^{(2)}\right)
\circ
\left(\bigoplus_r \Phi_n\right)
\circ
\Pi_k
\circ L_N^{(1)}.
\]
接下来，因为 \(\Phi_N=\phi^{\oplus N}\)、\(\Phi_n=\phi^{\oplus n}\)，而 \(\Pi_k\) 只是把坐标按 leaf layout 分组，所以逐坐标映射与该分组交换：
\[
\left(\bigoplus_r\Phi_n\right)\circ\Pi_k
=
\Pi_k\circ\Phi_N.
\]
为了看清这一点，取任意 \(x=(x_0,\ldots,x_{N-1})\)，第 \(r\) 个 leaf 的第 \(j\) 个坐标在左侧为 \(\phi(x_{r+jk})\)，在右侧也为 \(\phi(x_{r+jk})\)。因此上式成立。

于是整体变为
\[
\Merge_k\circ
\left(\bigoplus_r L_n^{(2)}\right)
\circ
\Pi_k
\circ
\Phi_N
\circ L_N^{(1)}.
\]
再由第二个兼容条件，\(\Merge_k\circ(\bigoplus_r L_n^{(2)})\circ\Pi_k=L_N^{(2)}\)。因此得到
\[
L_N^{(2)}\circ\Phi_N\circ L_N^{(1)}=F_N.
\]
定理得证。
\end{proof}

\begin{example}[CKKS core 是 sandwich theorem 的实例]
在 CKKS bootstrapping core 中，取 \(L^{(1)}=\CtwoS\)、\(\Phi=\EvalMod\)、\(L^{(2)}=\StwoC\)。Lemma~\ref{lem:split-c2s} 给出第一个线性兼容条件，Lemma~\ref{lem:s2c-merge} 给出第二个线性兼容条件，Lemma~\ref{lem:evalmod} 正是逐坐标中间层与 layout grouping 的交换性。因此 Theorem~\ref{thm:ckks-core-factorization} 是 Theorem~\ref{thm:general-sandwich} 的直接实例。
\end{example}

\subsection{本节小结}
\label{subsec:main-theory-summary}

本节证明了三件事。第一，\(\Split_k\) 后接 leaf \(\CtwoS_n\) 的效果等价于 top \(\CtwoS_N\) 后的 coefficient residue-class 分组。第二，\(\EvalMod\) 虽然是非线性的，但它是逐坐标函数，因此可以在分组后的 leaf vectors 上独立执行。第三，leaf \(\StwoC_n\) 后接 \(\Merge_k\) 可以重构 top \(\StwoC_N\) 的输出。三者合起来给出 CKKS bootstrapping core 的 subring factorization。

这也解释了为什么本文不需要在 split 后先做 slot-preserving recovery。本文分解的对象不是原始 big-ring canonical slots，而是 bootstrapping core 中 \(\CtwoS\) 后的 coefficient-packed coordinates。下一节将在该代数结论的基础上讨论 ciphertext-level algorithm、ring switching noise、复杂度和参数选择。

\todo{实现章节需要补充：Lattigo 的 half-slot convention、real/imag branches、DFT factorization、basis ordering 与本节 normalized layout 的对应关系。本节不展开 implementation scaling。}

\section{算法、正确性与参数分析}

\subsection{Factored bootstrapping algorithm}
Ciphertext-level algorithm 使用 ring switching 在真实密钥下实现 algebraic split。

\begin{algorithm}[h]
\caption{Subring-factored CKKS bootstrapping}
\begin{algorithmic}[1]
\Require 大环 ciphertext $ct_N$，ring-switching keys，leaf bootstrapping keys。
\State $(ct_{n,0},\ldots,ct_{n,k-1})\gets\RSDown_{N\to n}(ct_N)$。
\For{$r=0,\ldots,k-1$ 并行执行}
    \State $ct'_{n,r}\gets\BTS_n(ct_{n,r})$。
\EndFor
\State $ct'_N\gets\RSUp_{n\to N}(ct'_{n,0},\ldots,ct'_{n,k-1})$。
\State \Return $ct'_N$。
\end{algorithmic}
\end{algorithm}

代数定理说明 leaf-wise core 为什么正确。算法还要求 $\RSDown$ 输出 leaf secret 下有效的 leaf ciphertexts，并要求 $\RSUp$ 能回到目标 top secret。

\subsection{正确性与噪声}
假设 $ct_N$ 解密为 $m+e$。ring switching down 之后，我们要求
\[
\Dec_{s_n}(ct_{n,r})=m_r+e_r,
\]
其中 $m_r$ 是 $m$ 的第 $r$ 个 coefficient component，且 $\|e_r\|$ 在 leaf bootstrapping 的输入容忍范围内。如果每个 leaf bootstrapping 成功，则 $ct'_{n,r}$ 解密为 $m_r+e'_r$，其中 $e'_r$ 是刷新后的误差。merge 与 switch up 后得到一个解密为 $m+e_{\mathrm{out}}$ 的 ciphertext。

一个保守的 failure decomposition 是
\[
\Pr[\mathrm{fail}]\le \delta_{\mathrm{RSDown}}+k\delta_{\mathrm{BTS},n}+\delta_{\mathrm{RSUp}}.
\]
噪声来源包括输入 CKKS error、ring-switching key-switching error、leaf linear-transform error、EvalMod approximation error、rescale/RNS rounding error 和 final switch-up error。

\todo{给出更紧的噪声界，或者明确标注这一部分只是 conservative framework。}

\subsection{复杂度与 overhead}
令 $T_{\mathrm{shared}}(N)$ 表示仍保留在 top-level 的 preprocessing，令 $T_{\mathrm{core}}(N)$ 表示 $\CtwoS_N$、$\EvalMod_N$ 和 $\StwoC_N$ 的成本。Direct cost 建模为
\[
T_{\mathrm{dir}}(N)=T_{\mathrm{shared}}(N)+T_{\mathrm{core}}(N).
\]
Factored cost 建模为
\[
T_{\mathrm{fac}}(N,k)=T_{\mathrm{shared}}(N)+T_{\mathrm{RSDown}}(N,k)+T_{\mathrm{leaf}}^{\parallel}(n,k)+T_{\mathrm{RSUp}}(N,k).
\]
如果有 $w$ 个 workers，parallel efficiency 为 $\eta_{\mathrm{par}}$，则一阶近似为
$T_{\mathrm{leaf}}^{\parallel}\approx kT_{\mathrm{core}}(n)/(\min(k,w)\eta_{\mathrm{par}})$。预期加速来自用若干个小 core 替换一个大 core；额外成本来自 ring switching、extra switching keys、leaf bootstrapping keys、synchronization 和 memory。

\todo{补充 measured per-stage timings、key generation time、key sizes 和 peak memory。}

\subsection{安全性与参数选择}
安全性必须对所有出现公开 RLWE samples 的 rings 进行估计。特别地，leaf security 不会自动由 top-ring security 继承。如果 leaf ciphertexts 和 leaf evaluation keys 生活在 $R_{q,n}$ 中，则相关估计需要使用 leaf dimension $n$ 以及这些 keys 使用的最大公开 modulus，通常可保守近似为 $QP$。

这带来一个张力：更小的 $n$ 能加速 leaf computation，但 leaf modulus 必须足够小才能满足安全性。因此，一个 security-aware parameter table 至少应报告：top $N$ 与 $QP$、leaf $n$ 与 $Q_{\mathrm{leaf}}P_{\mathrm{leaf}}$、ring-switching key moduli、bootstrapping key moduli，以及所有估计中的最小安全位数。

\begin{table}[h]
\centering
\begin{tabular}{lllll}
\toprule
Tier & $\log N$ & $k$ & $\log n$ & 定位\\
\midrule
Prototype & 16 & 4 & 14 & 机制与正确性验证\\
Prototype+ & 16 & 2 & 15 & 正确性与 timing\\
Security-aware & 17 & 2 & 16 & 当前主实验\\
Heavy target & 18 & 4 & 16 & 后续高吞吐目标\\
\bottomrule
\end{tabular}
\caption{计划中的参数层级。}
\end{table}

\todo{对 top keys、leaf keys、RSDown keys、RSUp keys 和 bootstrapping keys 运行 lattice estimator。只有完成表格后，才能将 security-aware language 替换为具体 bit-security。}


\subsection{安全证明}

这里得加一个安全性证明。


\section{实现与实验评估}

\subsection{Implementation setup}
Prototype 基于 Lattigo 实现。实现使用 top parameter set 处理输入和输出 ciphertext，使用 leaf parameter sets 执行 subring bootstrappings，使用 ring-switching keys 实现 $\RSDown$ 和 $\RSUp$，并使用 leaf evaluators pool 执行并行计算。需要检查 top 与 leaf bases 的兼容性，包括 exact $Q/P$ basis、NTT-prime compatibility、scale conventions 和 layout ordering。

\todo{补充 Lattigo version、hardware、workers 数量、exact parameter literals 和 key-generation settings。}

\subsection{Layout and normalization diagnostics}
理想代数定理并不规定 Lattigo 的 implementation scaling。实践中我们观察到
\[
\BTS^{\mathrm{fac,impl}}_{N,k}\approx \gamma_{\mathrm{impl}}\BTS_N^{\mathrm{impl}}.
\]
其中 $\gamma_{\mathrm{impl}}$ 是 $\CtwoS$ 与 $\StwoC$ 的 effective scaling choices 带来的 implementation convention，不是 factorization theorem 的数学修正。有些参数组 RAW 形式直接对齐，另一些参数组需要显式除以一个小因子，例如 $k$。

因此，implementation 应记录 top $\CtwoS/\StwoC$ scaling、leaf $\CtwoS/\StwoC$ scaling、output scale、level、RAW error、normalized error 和 fitted $\gamma$。

\todo{插入 $\log N=17,k=2$ 和 $\log N=16,k=2$ 实验的 normalization diagnostic table。}

\subsection{Correctness experiments}
Core correctness experiments 应包括 Split/Merge roundtrip、$\CtwoS$ alignment、after-EvalMod alignment、final direct-vs-factored error、output level 和 normalization behavior。当前计划中的主实验是
\[
\log N=17,
\qquad k=2,
\qquad \log n=16.
\]
已有运行显示，该设置可以获得接近的 direct-vs-factored alignment 和显著 speedup。

\todo{插入最终表格，包括 C2S alignment、after-EvalMod alignment、final error、output level，以及多次 trials 下的 failure rate。}

\subsection{Performance and overhead}
Performance table 应报告 direct full bootstrapping time、factored full bootstrapping time、speedup、per-stage timings、worker count、key size 和 memory。更有用的表格应分别列出 $\RSDown$、leaf $\CtwoS$、leaf $\EvalMod$、leaf $\StwoC$ 和 $\RSUp$。

\todo{补 key size 和 peak memory。只有 timing 还不足以支撑完整 systems claim。}

\subsection{RNS-aware reduced-$Q$ investigation}
一个核心开放问题是：switch 到 subrings 之后，leaf bootstrapping 是否可以使用更小的 modulus chain $Q_{\mathrm{leaf}}$ 或更小的公开 modulus $Q_{\mathrm{leaf}}P_{\mathrm{leaf}}$。这很重要，因为 leaf security 要在 dimension $n$ 下估计。（待完成）

在 coefficient/RNS 层面，basis extension、modulus-chain restriction 和 coefficient-wise RNS operations 与 coefficient split 交换。然而，Lattigo 实验显示，现有 bootstrapping preprocessing，例如 $\mathsf{Evaluator.ModUp}$ 或组合的 $\mathsf{ScaleDown}+\mathsf{ModUp}$ block，不能简单地移到 split 之后再在 leaves 上独立调用。因此，正确的理论目标不是与 direct path 的中间态相等。Reduced-$Q$ factored bootstrapping 应该被表述为一个独立的 leaf bootstrapping design problem：low-modulus $\RSDown$ 必须输出有效 leaf ciphertexts，并且其噪声需要落在由 $Q_{\mathrm{leaf}}$ 参数化的 leaf bootstrapping 输入容忍范围内。

\todo{Reduced-$Q$ leaf bootstrapping 尚未完成。当前证据排除了最简单的 Lattigo preprocessing 重排，但没有排除 custom low-level preprocessing 或新的 leaf parameter set。}

\section{讨论与限制}

当前草稿已经建立 CKKS bootstrapping core 的代数分解、split 的 coefficient/slot interpretation，以及通过 subring ring switching 实现 ciphertext-level algorithmic path 的结构。它也给出了实现诊断和参数分析的框架。但仍有若干部分尚未完成。

\todo{Reduced-$Q$ leaf bootstrapping 仍然是开放问题。已有实验表明 Lattigo preprocessing 不能直接下放到 leaves；需要 custom RNS/preprocessing layer 或新的 leaf bootstrapping parameter design。}

\todo{完整安全估计缺失。最终论文必须对 top、leaf、RSDown、RSUp 和 bootstrapping keys 跑 lattice estimator，并报告最小 bit security。}

\todo{噪声分析目前是 conservative framework。需要加入更紧的 bound，或者至少加入经过验证的 empirical precision model。}

\todo{Key size、peak memory 和 per-stage timing 尚未完整。这些信息对于评估 ring switching 和 leaf keys 的 overhead 是必要的。}

\todo{更大 $k$ 的 security-aware 参数实验尚未完成。除非 leaf security 得到确认，否则 $k=4$ prototype 不应作为安全主实例呈现。}

最后，preliminaries 中讨论的 slot-lifting 或 slot-preserving transform 不是主算法的一部分。它是有用的比较对象，也有助于理解其他语义；但本文提出的 bootstrapping factorization 工作在 coefficient/module layout 中。

\section{结论}
本文草稿提出了一条 subring-factored CKKS bootstrapping 路线。主要结果是传统 $\CtwoS\to\EvalMod\to\StwoC$ core 与 coefficient/module decomposition 兼容：一个大环 core 可以被替换为独立的 leaf cores 加一次 merge。这不同于 ring packing、subring-secret acceleration 和 PaCo 这类新 bootstrapping 设计。本文保留传统 CKKS bootstrapping core，但改变其执行粒度。下一步需要完成严格证明、补齐实现表格、运行完整安全估计，并判断 reduced-$Q$ leaf bootstrapping 是否能够做到正确且安全。

\begin{thebibliography}{9}
\bibitem{ghs-ringswitch}
C. Gentry, S. Halevi, and N. P. Smart. Ring switching in BGV-style homomorphic encryption. Preliminary version, 2012.

\bibitem{hermes}
Y. Bae, J. H. Cheon, J. Kim, J. H. Park, and D. Stehle. HERMES: Efficient ring packing using MLWE ciphertexts and application to transciphering. 2024.

\bibitem{subring-secret}
S. Ma, T. Huang, A. Wang, and X. Wang. Practical dense-key bootstrapping with subring secret encapsulation. 2025.

\bibitem{paco}
J.-S. Coron and T. Seure. PaCo: Bootstrapping for CKKS via partial CoeffToSlot. 2025.

\bibitem{ckks}
J. H. Cheon, A. Kim, M. Kim, and Y. Song. Homomorphic encryption for arithmetic of approximate numbers. ASIACRYPT 2017.
\end{thebibliography}

\end{document}
